\documentclass[aspectratio=169,11pt]{beamer}

% ==================== 基础包 ====================
\usepackage[UTF8]{ctex}
\usepackage{amsmath,amssymb,amsthm}
\usepackage{mathtools}
\usepackage{bm}
\usepackage{graphicx}
\usepackage{booktabs}
\usepackage{array}
\usepackage{multirow}
\usepackage{xcolor}
\usepackage{listings}
\usepackage{tikz}
\usepackage{etoolbox}
\usetikzlibrary{shapes,arrows,positioning,calc}

% ==================== 颜色定义 ====================
\definecolor{darkblue}{RGB}{26,82,118}
\definecolor{lightblue}{RGB}{214,234,248}
\definecolor{accentgreen}{RGB}{39,174,96}
\definecolor{warnred}{RGB}{192,57,43}
\definecolor{codebg}{RGB}{248,248,248}
\definecolor{codeframe}{RGB}{200,200,200}
\definecolor{darkgray}{RGB}{60,60,60}

% ==================== Beamer主题设置 ====================
\usetheme{Madrid}
\usecolortheme[named=darkblue]{structure}
\setbeamercolor{title}{fg=white,bg=darkblue}
\setbeamercolor{frametitle}{fg=white,bg=darkblue}
\setbeamercolor{block title}{fg=white,bg=darkblue}
\setbeamercolor{block body}{bg=lightblue}
\setbeamercolor{block title example}{fg=white,bg=accentgreen}
\setbeamercolor{block body example}{bg=green!5}
\setbeamercolor{block title alerted}{fg=white,bg=warnred}
\setbeamercolor{block body alerted}{bg=red!5}
\setbeamercolor{item}{fg=darkblue}
\setbeamercolor{subitem}{fg=darkblue}
\setbeamercolor{itemize item}{fg=darkblue}
\setbeamercolor{itemize subitem}{fg=darkblue}
\setbeamercolor{enumerate item}{fg=darkblue}
\setbeamercolor{enumerate subitem}{fg=darkblue}

% ---- 修复Headline：只显示当前section名，不显式导航条 ----
\setbeamertemplate{headline}{%
  \leavevmode%
  \hbox{%
    \begin{beamercolorbox}[wd=\paperwidth,ht=2.2ex,dp=0.8ex,leftskip=10pt]{palette quaternary}%
      \usebeamerfont{section in head/foot}\insertsectionhead%
    \end{beamercolorbox}%
  }%
}

% ---- 修复Footline：确保三部分宽度正确 ----
\setbeamertemplate{footline}{%
  \leavevmode%
  \hbox{%
    \begin{beamercolorbox}[wd=.25\paperwidth,ht=2.25ex,dp=1ex,center]{author in head/foot}%
      \usebeamerfont{author in head/foot}\insertshortauthor%
    \end{beamercolorbox}%
    \begin{beamercolorbox}[wd=.50\paperwidth,ht=2.25ex,dp=1ex,center]{title in head/foot}%
      \usebeamerfont{title in head/foot}\insertshorttitle%
    \end{beamercolorbox}%
    \begin{beamercolorbox}[wd=.25\paperwidth,ht=2.25ex,dp=1ex,right]{date in head/foot}%
      \usebeamerfont{date in head/foot}\insertframenumber{} / \inserttotalframenumber\hspace*{2ex}%
    \end{beamercolorbox}%
  }%
}

% ==================== 代码样式 ====================
\lstset{
    language=R,
    basicstyle=\ttfamily\footnotesize,
    backgroundcolor=\color{codebg},
    frame=single,
    rulecolor=\color{codeframe},
    breaklines=true,
    breakatwhitespace=true,
    showstringspaces=false,
    keywordstyle=\color{darkblue}\bfseries,
    commentstyle=\color{gray}\itshape,
    stringstyle=\color{accentgreen},
    numbers=left,
    numberstyle=\tiny\color{gray},
    numbersep=5pt,
    tabsize=2,
    xleftmargin=10pt,
    xrightmargin=5pt,
    aboveskip=5pt,
    belowskip=5pt,
    extendedchars=true,
    columns=fullflexible,
    keepspaces=true,
    morekeywords={mmrm,us,cs,ar1,toep,emmeans,pairs,VarCorr,summary,library,data,proc,mixed,class,model,repeated,lsmeans,ddfm,method,reml,fit,formula,control,optimizer,vcov,tidy,glance,augment,component,df_1d,df_md,contrast}
}

% ==================== 数学符号简写 ====================
\newcommand{\bY}{\mathbf{Y}}
\newcommand{\by}{\mathbf{y}}
\newcommand{\bX}{\mathbf{X}}
\newcommand{\bZ}{\mathbf{Z}}
\newcommand{\bbeta}{\boldsymbol{\beta}}
\newcommand{\btheta}{\boldsymbol{\theta}}
\newcommand{\balpha}{\boldsymbol{\alpha}}
\newcommand{\bepsilon}{\boldsymbol{\epsilon}}
\newcommand{\bSigma}{\boldsymbol{\Sigma}}
\newcommand{\bOmega}{\boldsymbol{\Omega}}
\newcommand{\bPsi}{\boldsymbol{\Psi}}
\newcommand{\bLambda}{\boldsymbol{\Lambda}}
\newcommand{\bI}{\mathbf{I}}
\newcommand{\bD}{\mathbf{D}}
\newcommand{\bP}{\mathbf{P}}
\newcommand{\bS}{\mathbf{S}}
\newcommand{\bi}{\mathbf{b}_i}
\newcommand{\Var}{\mathrm{Var}}
\newcommand{\Cov}{\mathrm{Cov}}
\newcommand{\E}{\mathrm{E}}
\newcommand{\tr}{\mathrm{tr}}
\newcommand{\diag}{\mathrm{diag}}
\newcommand{\argmin}{\mathop{\arg\min}}

% ==================== 标题信息 ====================
\title[MMRM教学]{Mixed Model for Repeated Measures (MMRM)\\
LMM与MMRM}
\author[MMRM Tutorial]{Guo Wei}
\institute{Department of Health Statistics, NMU}
\date{2026-06-11}

% ==================== 文档开始 ====================
\begin{document}

%---------------------------------------------------
\begin{frame}[plain]
\titlepage
\end{frame}

%---------------------------------------------------
\begin{frame}{目录}
\tableofcontents
\end{frame}

%===================================================
\section{引言：纵向数据与重复测量}
%===================================================

%---------------------------------------------------
\begin{frame}{纵向数据（Longitudinal Data）}
\begin{block}{定义}
对同一受试者（subject）在\textbf{多个时间点}重复测量同一指标所得到的数据。
\end{block}

\vspace{0.2cm}
\textbf{典型场景：}
\begin{itemize}
    \item 临床试验：基线、第2周、第4周、第8周的疗效评分
    \item 流行病学：随访调查中的多次血压测量
    \item 心理学：干预前、干预后、随访期的量表得分
\end{itemize}

\vspace{0.2cm}
\begin{alertblock}{核心挑战}
\begin{enumerate}
    \item \textbf{相关性}：同一受试者的多次测量之间存在相关性
    \item \textbf{缺失数据}：受试者可能在某些时间点失访或退出
    \item \textbf{异方差}：不同时间点的变异程度可能不同
\end{enumerate}
\end{alertblock}
\end{frame}

%---------------------------------------------------
\begin{frame}{分析策略的演进}
\begin{center}
\begin{tikzpicture}[
    node distance=1.4cm,
    box/.style={rectangle, draw=darkblue, fill=lightblue, rounded corners,
                minimum width=2.4cm, minimum height=1cm, align=center, font=\footnotesize},
    arrow/.style={->, >=stealth, thick, darkblue}
]
\node[box] (ancova) {ANCOVA\\（单时间点）};
\node[box, right=of ancova] (locf) {LOCF\\（末次观测结转）};
\node[box, right=of locf] (gee) {GEE\\（广义估计方程）};
\node[box, right=of gee] (mmrm) {MMRM\\（重复测量混合）};

\draw[arrow] (ancova) -- (locf);
\draw[arrow] (locf) -- (gee);
\draw[arrow] (gee) -- (mmrm);

\node[below=0.25cm of locf, font=\tiny, text=warnred] {有偏估计，低估SE};
\node[below=0.25cm of gee, font=\tiny, text=darkblue] {稳健但效率低};
\node[below=0.25cm of mmrm, font=\tiny, text=accentgreen] {\textbf{当前金标准}};
\end{tikzpicture}
\end{center}

\vspace{0.3cm}
\begin{block}{为什么MMRM成为金标准？}
\begin{itemize}
    \item 2010年NRC报告明确反对LOCF等单一插补法
    \item EMA缺失数据指南推荐基于似然的方法
    \item MMRM在MAR假设下无需显式插补即可利用所有观测数据
\end{itemize}
\end{block}
\end{frame}

%===================================================
\section{传统线性混合模型（LMM）原理}
%===================================================

%---------------------------------------------------
\begin{frame}{LMM的基本框架}
\begin{block}{Laird \& Ware (1982) 基本线性混合效应模型}
对第 $i$ 个受试者，其 $n_i$ 维响应向量 $\by_i$：
\[
\by_i = \bX_i \bbeta + \bZ_i \bi + \bepsilon_i, \quad i = 1, \dots, M
\]
\end{block}

\vspace{0.1cm}
\begin{columns}[T,totalwidth=\textwidth]
\begin{column}{0.48\textwidth}
\textbf{模型组成：}
\begin{itemize}
    \item $\bbeta$：$p$ 维\textbf{固定效应}向量
    \item $\bi$：$q$ 维\textbf{随机效应}向量
    \item $\bX_i$：$n_i \times p$ 固定效应设计矩阵
    \item $\bZ_i$：$n_i \times q$ 随机效应设计矩阵
\end{itemize}
\end{column}
\begin{column}{0.48\textwidth}
\textbf{分布假设：}
\begin{align*}
\bi &\sim N(\mathbf{0}, \bPsi) \\
\bepsilon_i &\sim N(\mathbf{0}, \sigma^2 \bI) \\
\bi &\perp\!\!\!\perp \bepsilon_i
\end{align*}
\end{column}
\end{columns}

\vspace{0.1cm}
\begin{exampleblock}{关键特征}
随机效应 $\bi$ 描述每个受试者相对于总体均值的\textbf{偏移}（均值为0）。
\end{exampleblock}
\end{frame}

%---------------------------------------------------
\begin{frame}{LMM的方差-协方差分解}
\begin{block}{响应变量的方差-协方差矩阵}
\[
\Var(\by_i) = \bSigma_i = \underbrace{\bZ_i \bPsi \bZ_i^\top}_{\text{G-side（随机效应）}} + \underbrace{\sigma^2 \bI}_{\text{R-side（残差）}}
\]
\end{block}

\vspace{0.2cm}
\textbf{两部分竞争与权衡：}
\begin{itemize}
    \item \textbf{G-side}：通过随机效应捕捉受试者\textbf{间}变异
    \item \textbf{R-side}：通过残差结构捕捉受试者\textbf{内}相关
\end{itemize}

\vspace{0.2cm}
\begin{alertblock}{重要提醒}
当两部分同时复杂时，可能出现\textbf{不可识别性}问题。实践中通常只使用其中一侧。
\end{alertblock}
\end{frame}

%---------------------------------------------------
\begin{frame}{LMM的扩展：异方差与相关残差}
\begin{block}{扩展模型}
将基本LMM的残差假设放宽为：$\bepsilon_i \sim N(\mathbf{0}, \bLambda_i)$，其中 $\bLambda_i$ 由参数 $\btheta$ 参数化。
\end{block}

\textbf{扩展后的完整方差-协方差：}\vspace{-0.2cm}
\[
\Var(\by_i) = \bSigma_i = \bZ_i \bPsi \bZ_i^\top + \bLambda_i
\]

\vspace{-0.1cm}
\begin{columns}[T,totalwidth=\textwidth]
\begin{column}{0.48\textwidth}
\textbf{可建模的R-side结构：}
\begin{itemize}
    \item 异方差（不同时间点方差不同）
    \item 自相关（AR1）、复合对称（CS）
    \item 非结构化（UN）、Toeplitz
\end{itemize}
\end{column}
\begin{column}{0.48\textwidth}
\begin{exampleblock}{文献记号}
在SAS和混合模型文献中：
\begin{itemize}
    \item G-side = 随机效应成分
    \item R-side = 残差/within-subject成分
\end{itemize}
（Cnaan, Laird \& Slasor, 1997）
\end{exampleblock}
\end{column}
\end{columns}
\end{frame}

%---------------------------------------------------
\begin{frame}{LMM在纵向数据中的局限}
\begin{block}{基本LMM的隐含假设}
原始LMM假设 within-subject 误差满足 $\bepsilon_i \sim N(\mathbf{0}, \sigma^2 \bI)$，即\textbf{独立、同方差}。
\end{block}

\vspace{0.2cm}
\textbf{纵向数据中的现实：}
\begin{enumerate}
    \item \textbf{相关性}：同一受试者相邻时间点测量高度相关
    \item \textbf{异方差}：后期测量的变异通常大于早期
    \item \textbf{不平衡}：不同受试者观测时间点可能不同
\end{enumerate}

\vspace{0.2cm}
\begin{alertblock}{结论}
仅通过G-side随机效应往往\textbf{不足以}充分刻画within-subject相关结构。需要更灵活的R-side建模——这正是MMRM的核心思想。
\end{alertblock}
\end{frame}

%---------------------------------------------------
\begin{frame}{LMM的估计：REML vs ML}
\begin{columns}[T,totalwidth=\textwidth]
\begin{column}{0.48\textwidth}
\begin{block}{两种似然估计}
\begin{itemize}
    \item \textbf{ML}：最大化完整似然
    \item \textbf{REML}：最大化残差似然，对固定效应校正
\end{itemize}
\end{block}

\textbf{REML的优势：}
\begin{itemize}
    \item 方差分量估计\textbf{偏差更小}
    \item 平衡设计中等价于ANOVA估计
    \item 临床试验\textbf{默认使用REML}
\end{itemize}

\textbf{ML的用途：}
\begin{itemize}
    \item 模型比较（AIC、BIC）
    \item 嵌套模型似然比检验
\end{itemize}
\end{column}
\begin{column}{0.48\textwidth}
\begin{exampleblock}{REML目标函数}
\[
\ell_{REML}(\btheta) = \ell_{ML}(\btheta) - \frac{1}{2} \log |\bX^\top \bSigma^{-1} \bX|
\]
\vspace{0.1cm}
第二项为固定效应的校正项。
\end{exampleblock}
\end{column}
\end{columns}
\end{frame}

%===================================================
\section{MMRM原理与模型结构}
%===================================================

%---------------------------------------------------
\begin{frame}{MMRM的定义}
\begin{block}{Mixed Model for Repeated Measures (MMRM)}
MMRM是一种基于似然的方法，用于分析在\textbf{同一受试者上重复测量}的连续结局变量。核心特征：
\begin{enumerate}
    \item \textbf{仅含固定效应}的均值模型
    \item 通过\textbf{结构化残差协方差矩阵} $\bSigma$ 直接建模受试者内相关性
    \item 参数通过\textbf{REML}估计
\end{enumerate}
\end{block}

\vspace{0.2cm}
\begin{alertblock}{重要澄清}
MMRM虽然名为``混合模型''，但经典临床试验中的MMRM\textbf{通常不含随机效应}（如随机截距或斜率）。焦点在\textbf{残差协方差结构}（R-side）而非随机效应（G-side）。
\end{alertblock}
\end{frame}

%---------------------------------------------------
\begin{frame}{MMRM的数学模型}
\begin{block}{个体层面模型}
对第 $i$ 个受试者，其基线后测量向量 $\by_i$：
\[
\by_i = \bX_i \bbeta + \bepsilon_i, \quad \bepsilon_i \sim N(\mathbf{0}, \bSigma_i)
\]
\end{block}

\vspace{-0.1cm}
\textbf{其中：}
\begin{itemize}
    \item $\bX_i$：固定效应设计矩阵（治疗、访视、交互、协变量）
    \item $\bbeta$：固定效应回归系数
    \item $\bSigma_i$：受试者 $i$ 的within-subject协方差矩阵
\end{itemize}

\vspace{0.1cm}
\begin{block}{试验整体模型}
将所有 $n$ 个受试者合并：$\bY = \bX \bbeta + \bepsilon$，\quad $\bepsilon \sim N(\mathbf{0}, \bOmega)$

其中 $\bOmega = \diag(\bSigma_1, \bSigma_2, \dots, \bSigma_M)$ 为\textbf{块对角矩阵}。
\end{block}
\end{frame}

%---------------------------------------------------
\begin{frame}{MMRM作为LMM的特例}
\begin{block}{从LMM到MMRM的简化}
在LMM的一般形式中：$\Var(\by_i) = \underbrace{\bZ_i \bPsi \bZ_i^\top}_{\text{G-side}} + \underbrace{\bLambda_i}_{\text{R-side}}$

MMRM选择\textbf{省略G-side}，仅保留R-side：$\Var(\by_i) = \bSigma_i = \bLambda_i$
\end{block}

\vspace{0.2cm}
\begin{columns}[T,totalwidth=\textwidth]
\begin{column}{0.52\textwidth}
\textbf{为什么可以省略G-side？}
\begin{itemize}
    \item 纵向试验中within-subject相关性是主要关注点
    \item 结构化协方差矩阵已能充分捕捉相关性
    \item 避免G-side与R-side的不可识别问题
\end{itemize}
\end{column}
\begin{column}{0.44\textwidth}
\begin{exampleblock}{\footnotesize Mallinckrodt et al. (2008)}
\footnotesize ``在仅有一级分组的纵向研究中，within-subject成分尤为重要，额外的随机效应成分往往并非严格必需。''
\end{exampleblock}
\end{column}
\end{columns}
\end{frame}

%---------------------------------------------------
\begin{frame}{MMRM的固定效应结构}
\begin{block}{典型的确证性试验MMRM}
$\bY = \bX \bbeta + \bepsilon$
\end{block}

\vspace{0.1cm}
\textbf{固定效应通常包括：}
\begin{enumerate}
    \item \textbf{治疗组}（Treatment）：随机化分组
    \item \textbf{访视}（Visit）：作为\textbf{分类因子}，每个计划时间点单独估计
    \item \textbf{治疗$\times$访视交互}：得到每个访视的治疗效应
    \item \textbf{基线值}（Baseline）：结局变量的基线测量
    \item \textbf{基线$\times$访视交互}（可选）
    \item \textbf{分层因子}：随机化分层因素
\end{enumerate}

\vspace{0.1cm}
\begin{alertblock}{关键设计选择}
将访视作为\textbf{分类变量}（而非连续时间）是MMRM的标志性特征——允许治疗效应随时间呈现\textbf{任意形状}。
\end{alertblock}
\end{frame}

%---------------------------------------------------
\begin{frame}{MMRM中的协方差矩阵构造}
\begin{block}{个体协方差矩阵}
设共有 $m$ 个计划访视，总体协方差矩阵为 $\bSigma \in \mathbb{R}^{m \times m}$。对第 $i$ 个受试者：
\vspace{-0.2cm}
\[
\bSigma_i = \bS_i^\top \bSigma \bS_i
\]
\vspace{-0.2cm}
其中 $\bS_i \in \{0,1\}^{m \times m_i}$ 为受试者特定的``子集矩阵''。
\end{block}

\textbf{示例}（4个访视，受试者缺失VIS3）：
\[
\small
\bS_i = \begin{pmatrix} 1 & 0 & 0 \\ 0 & 1 & 0 \\ 0 & 0 & 0 \\ 0 & 0 & 1 \end{pmatrix}, \quad
\bSigma_i = \begin{pmatrix} \sigma_{11} & \sigma_{12} & \sigma_{14} \\ \sigma_{21} & \sigma_{22} & \sigma_{24} \\ \sigma_{41} & \sigma_{42} & \sigma_{44} \end{pmatrix}
\]

\begin{exampleblock}{块对角整体矩阵}
\vspace{-0.2cm}
$\bOmega = \diag(\bSigma_1, \bSigma_2, \dots, \bSigma_M)$\quad\quad 受试者间独立，受试者内相关。
\end{exampleblock}
\end{frame}

%---------------------------------------------------
\begin{frame}{MMRM的估计与推断}
\begin{block}{REML估计}
\vspace{-0.2cm}
MMRM通过最大化残差似然估计 $\bbeta$ 和协方差参数 $\btheta$：
\[
\small
\ell_{REML}(\btheta) = -\frac{1}{2} \sum_{i=1}^{M} \left[ \log |\bSigma_i| + (\by_i - \bX_i \hat{\bbeta})^\top \bSigma_i^{-1} (\by_i - \bX_i \hat{\bbeta}) \right] + C
\]
\end{block}

\vspace{-0.3cm}
\begin{columns}[T,totalwidth=\textwidth]
\begin{column}{0.58\textwidth}
\textbf{固定效应估计（GLS）：}
\vspace{-0.3cm}
\[
\footnotesize
\hat{\bbeta} = \left( \sum_{i} \bX_i^\top \bSigma_i^{-1} \bX_i \right)^{-1} \sum_{i} \bX_i^\top \bSigma_i^{-1} \by_i
\]
\end{column}
\begin{column}{0.38\textwidth}
\textbf{方差-协方差：}
\vspace{-0.3cm}
\[
\footnotesize
\Var(\hat{\bbeta}) = \left( \sum_{i} \bX_i^\top \bSigma_i^{-1} \bX_i \right)^{-1}
\]
\end{column}
\end{columns}
\end{frame}

%===================================================
\section{协方差结构详解}
%===================================================

%---------------------------------------------------
\begin{frame}{协方差结构概述}
\begin{block}{协方差矩阵的分解}
多数协方差结构采用：$\bSigma = \bD \bP \bD$

其中 $\bD = \diag(\sigma_1, \dots, \sigma_m)$ 为标准差对角矩阵，$\bP$ 为相关矩阵。
\end{block}

\textbf{``异质性''（Heterogeneous）vs ``同质性''（Homogeneous）：}
\begin{itemize}
    \item \textbf{异质性}：允许不同时间点有不同方差
    \item \textbf{同质性}：假设所有时间点方差相同
\end{itemize}

\begin{alertblock}{参数数量}
参数数量 $k$ 取决于结构复杂度。UN需要 $k = m(m+1)/2$ 个参数，是``最自由''的结构。
\end{alertblock}
\end{frame}

%---------------------------------------------------
\begin{frame}{非结构化协方差（Unstructured, UN）}
\begin{block}{定义}
不对协方差矩阵施加任何结构约束——每个方差和每个协方差都单独估计。
\end{block}

\vspace{0.1cm}
\textbf{4个访视的UN矩阵：}
\vspace{-0.2cm}
\[
\bSigma = \begin{pmatrix}
\sigma_1^2 & \sigma_{12} & \sigma_{13} & \sigma_{14} \\
\sigma_{21} & \sigma_2^2 & \sigma_{23} & \sigma_{24} \\
\sigma_{31} & \sigma_{32} & \sigma_3^2 & \sigma_{34} \\
\sigma_{41} & \sigma_{42} & \sigma_{43} & \sigma_4^2
\end{pmatrix}
\]

\vspace{-0.1cm}
\begin{columns}[T,totalwidth=\textwidth]
\begin{column}{0.48\textwidth}
\textbf{优点：}
\begin{itemize}
    \item 假设最少，最灵活
    \item 确证性试验的\textbf{监管默认}
\end{itemize}
\end{column}
\begin{column}{0.48\textwidth}
\textbf{缺点：}
\begin{itemize}
    \item 参数最多：$k = m(m+1)/2$
    \item 访视多或样本小时\textbf{可能不收敛}
\end{itemize}
\end{column}
\end{columns}
\end{frame}

%---------------------------------------------------
\begin{frame}{复合对称（CS）与一阶自回归（AR1）}
\begin{columns}[T,totalwidth=\textwidth]
\begin{column}{0.48\textwidth}
\begin{block}{\footnotesize 复合对称 CS}
\footnotesize
假设任意两个时间点间相关恒定：$\rho_{ij} = \rho$

\textbf{同质CS}（$k=2$）：
\vspace{-0.2cm}
\[
\footnotesize
\bSigma = \sigma^2 \begin{pmatrix}
1 & \rho & \rho & \rho \\
\rho & 1 & \rho & \rho \\
\rho & \rho & 1 & \rho \\
\rho & \rho & \rho & 1
\end{pmatrix}
\]
\vspace{-0.1cm}
\textcolor{warnred}{\footnotesize\textbf{注意}：纵向数据中CS通常不现实}
\end{block}
\end{column}
\begin{column}{0.48\textwidth}
\begin{block}{\footnotesize 一阶自回归 AR1}
\footnotesize
相关随时间间隔几何衰减：$\rho_{ij} = \rho^{|i-j|}$

\textbf{同质AR1}（$k=2$）：
\vspace{-0.2cm}
\[
\footnotesize
\bSigma = \sigma^2 \begin{pmatrix}
1 & \rho & \rho^2 & \rho^3 \\
\rho & 1 & \rho & \rho^2 \\
\rho^2 & \rho & 1 & \rho \\
\rho^3 & \rho^2 & \rho & 1
\end{pmatrix}
\]
\vspace{-0.1cm}
\textcolor{accentgreen}{\footnotesize\textbf{适用}：等距访视，相关随时间递减}
\end{block}
\end{column}
\end{columns}
\end{frame}

%---------------------------------------------------
\begin{frame}{Toeplitz与Ante-dependence}
\begin{columns}[T,totalwidth=\textwidth]
\begin{column}{0.48\textwidth}
\begin{block}{\footnotesize Toeplitz}
\footnotesize
相关仅取决于时间\textbf{间隔}（lag）：$\rho_{ij} = \rho_{|i-j|}$

\textbf{4个访视的Toeplitz}：
\vspace{-0.2cm}
\[
\footnotesize
\bP = \begin{pmatrix}
1 & \rho_1 & \rho_2 & \rho_3 \\
\rho_1 & 1 & \rho_1 & \rho_2 \\
\rho_2 & \rho_1 & 1 & \rho_1 \\
\rho_3 & \rho_2 & \rho_1 & 1
\end{pmatrix}
\]
\vspace{-0.1cm}
参数：$k = m$（同质）或 $2m-1$（异质）
\end{block}
\end{column}
\begin{column}{0.48\textwidth}
\begin{block}{\footnotesize Ante-dependence (AD)}
\footnotesize
一阶ante-dependence模型：$\rho_{ij} = \prod_{k=i}^{j-1} \rho_k$

\textbf{4个访视的AD}：
\vspace{-0.2cm}
\[
\footnotesize
\bP = \begin{pmatrix}
1 & \rho_1 & \rho_1\rho_2 & \rho_1\rho_2\rho_3 \\
\rho_1 & 1 & \rho_2 & \rho_2\rho_3 \\
\rho_1\rho_2 & \rho_2 & 1 & \rho_3 \\
\rho_1\rho_2\rho_3 & \rho_2\rho_3 & \rho_3 & 1
\end{pmatrix}
\]
\vspace{-0.1cm}
参数：$k = m$（同质）或 $2m-1$（异质）
\end{block}
\end{column}
\end{columns}
\end{frame}

%---------------------------------------------------
\begin{frame}{协方差结构选择策略}
\begin{block}{确证性试验的推荐策略（Mallinckrodt et al.）}
\begin{enumerate}
    \item \textbf{首选UN}：预设非结构化协方差为首要分析
    \item \textbf{预设fallback}：在方案中预先指定UN不收敛时的替代结构
    \item \textbf{常见fallback}：Toeplitz或AR1（异质性版本）
    \item \textbf{避免CS}：复合对称在纵向数据中通常不现实
\end{enumerate}
\end{block}

\vspace{0.2cm}
\begin{center}
\small
\begin{tabular}{lccc}
\toprule
\textbf{结构} & \textbf{参数} & \textbf{假设强度} & \textbf{适用场景} \\
\midrule
UN & $m(m+1)/2$ & 最弱 & 小样本、确证性试验 \\
Toeplitz & $2m-1$ & 中等 & 相关仅依赖lag \\
AR1 & $m+1$ & 较强 & 等距访视、几何衰减 \\
CS & $m+1$ & 最强 & 极少用于纵向数据 \\
\bottomrule
\end{tabular}
\end{center}
\end{frame}

%===================================================
\section{缺失数据机制与MMRM的优势}
%===================================================

%---------------------------------------------------
\begin{frame}{缺失数据的三类机制}
\begin{block}{Rubin (1976) 缺失数据分类}
\begin{enumerate}
    \item \textbf{MCAR}（Missing Completely at Random）
    \begin{itemize}
        \item 缺失与任何数据（观测或未观测）均无关
        \item 例如：实验室设备故障导致某批次样本丢失
    \end{itemize}
    \item \textbf{MAR}（Missing at Random）
    \begin{itemize}
        \item 缺失仅依赖于\textbf{观测数据}（如早期测量值、治疗组、基线）
        \item 例如：病情较重的患者更可能退出
    \end{itemize}
    \item \textbf{MNAR}（Missing Not at Random）
    \begin{itemize}
        \item 缺失依赖于\textbf{未观测值本身}
        \item 例如：患者因疗效差而退出，且该疗效差未被记录
    \end{itemize}
\end{enumerate}
\end{block}
\end{frame}

%---------------------------------------------------
\begin{frame}{MMRM在MAR下的有效性}
\begin{block}{MMRM的核心优势}
作为\textbf{基于似然的方法}，MMRM在MAR假设下：
\begin{itemize}
    \item 无需任何显式插补
    \item 自动利用每个受试者的\textbf{所有可用观测}
    \item 通过联合建模所有访视，隐式处理缺失
\end{itemize}
\end{block}

\vspace{0.1cm}
\textbf{似然构建方式：}对第 $i$ 个受试者，仅使用其观测访视构建似然：
\vspace{-0.2cm}
\[
\small
\ell_i(\bbeta, \btheta) \propto -\frac{1}{2} \left[ \log |\bSigma_i^{obs}| + (\by_i^{obs} - \bX_i^{obs}\bbeta)^\top (\bSigma_i^{obs})^{-1} (\by_i^{obs} - \bX_i^{obs}\bbeta) \right]
\]

\vspace{0.1cm}
\begin{exampleblock}{与LOCF的对比}
\begin{itemize}
    \item LOCF：将末次观测``结转''到终点，\textbf{低估变异}
    \item 完整案例分析：丢弃所有缺失受试者，\textbf{损失信息}
    \item MMRM：利用所有信息，在MAR下\textbf{无偏且有效}
\end{itemize}
\end{exampleblock}
\end{frame}

%---------------------------------------------------
\begin{frame}{方法对比总结}
\begin{center}
\small
\begin{tabular}{lccc}
\toprule
\textbf{方法} & \textbf{使用数据} & \textbf{缺失处理} & \textbf{有效性} \\
\midrule
LOCF + ANCOVA & 单终点 & 单一插补 & \textcolor{warnred}{有偏，低估SE} \\
完整案例ANCOVA & 完成者 & 删除缺失 & \textcolor{warnred}{损失功效} \\
GEE & 所有观测 & 工作相关矩阵 & 需MCAR才无偏 \\
\textbf{MMRM} & \textbf{所有访视联合} & \textbf{似然隐式处理} & \textcolor{accentgreen}{\textbf{MAR下有效}} \\
多重插补 & 所有观测 & 多次插补 & MAR下有效 \\
\bottomrule
\end{tabular}
\end{center}

\vspace{0.2cm}
\begin{alertblock}{监管立场}
\begin{itemize}
    \item 2010年NRC报告：明确反对LOCF等单一插补法作为主要分析
    \item EMA缺失数据指南：推荐基于似然的方法（MMRM）和多重插补
    \item FDA/EMA审评中，MMRM是连续终点确证性试验的\textbf{de facto标准}
\end{itemize}
\end{alertblock}
\end{frame}

%===================================================
\section{LMM与MMRM：联系、区别与选择}
%===================================================

%---------------------------------------------------
\begin{frame}{LMM与MMRM的联系}
\begin{block}{数学层面的包含关系}
MMRM是LMM的\textbf{特例}：
\begin{itemize}
    \item LMM的一般形式：$\by_i = \bX_i\bbeta + \bZ_i\bi + \bepsilon_i$
    \item MMRM的形式：$\by_i = \bX_i\bbeta + \bepsilon_i$，其中 $\bepsilon_i \sim N(\mathbf{0}, \bSigma_i)$
\end{itemize}
即MMRM = LMM中令 $\bZ_i\bi = \mathbf{0}$，并将所有相关性放入R-side。
\end{block}

\vspace{0.2cm}
\textbf{共同特征：}
\begin{enumerate}
    \item 均基于线性模型框架，均使用REML/ML进行估计
    \item 均可处理不平衡数据，均支持协变量调整
    \item 在SAS中均通过PROC MIXED实现
\end{enumerate}
\end{frame}

%---------------------------------------------------
\begin{frame}{LMM与MMRM的核心区别}
\begin{center}
\small
\begin{tabular}{p{3cm}p{4.2cm}p{4.2cm}}
\toprule
\textbf{特征} & \textbf{传统LMM} & \textbf{MMRM} \\
\midrule
随机效应 & 通常包含（随机截距/斜率） & \textbf{通常不含} \\
相关建模 & G-side + R-side & \textbf{仅R-side} \\
协方差焦点 & 受试者间变异 & 受试者内相关 \\
典型应用 & 多中心试验、聚类数据 & 纵向重复测量 \\
参数解释 & 随机效应方差有意义 & 协方差结构为``讨厌参数'' \\
SAS语句 & \texttt{random} + \texttt{repeated} & 主要\texttt{repeated} \\
\bottomrule
\end{tabular}
\end{center}

\vspace{0.2cm}
\begin{exampleblock}{一句话总结}
\begin{itemize}
    \item LMM：``我关心受试者之间的差异有多大''
    \item MMRM：``我关心同一受试者多次测量如何相关，但只想得到治疗效应的准确估计''
\end{itemize}
\end{exampleblock}
\end{frame}

%---------------------------------------------------
\begin{frame}{应用场景选择指南}
\begin{columns}[T,totalwidth=\textwidth]
\begin{column}{0.48\textwidth}
\begin{block}{选择LMM当…}
\begin{itemize}
    \item 研究问题关注\textbf{受试者间变异}
    \item 需要预测个体轨迹
    \item 多层级结构（如患者嵌套在医院内）
    \item 随机效应本身有科学意义
\end{itemize}
\end{block}
\end{column}
\begin{column}{0.48\textwidth}
\begin{block}{选择MMRM当…}
\begin{itemize}
    \item 主要目标是\textbf{组间比较}（如治疗vs安慰剂）
    \item 数据来自临床试验的\textbf{计划访视}
    \item 需要处理\textbf{MAR缺失数据}
    \item 需要与SAS/监管标准保持一致
\end{itemize}
\end{block}
\end{column}
\end{columns}

\vspace{0.2cm}
\begin{alertblock}{注意}
在某些场景下两者可\textbf{混合使用}：例如多中心纵向试验中，中心作为随机效应（G-side），访视内相关用R-side结构建模。
\end{alertblock}
\end{frame}

%===================================================
\section{R实现：mmrm包与nlme}
%===================================================

%---------------------------------------------------
\begin{frame}[fragile]{mmrm包简介}
\begin{columns}[T,totalwidth=\textwidth]
\begin{column}{0.50\textwidth}
\begin{block}{openpharma/mmrm}
\footnotesize
由openpharma社区维护的R包，专为临床试验MMRM设计：
\begin{itemize}
    \item 快速、稳定的优化器
    \item 支持多种协方差结构
    \item 内置Satterthwaite和Kenward-Roger自由度调整
    \item 与\texttt{emmeans}无缝集成
\end{itemize}
\end{block}
\end{column}
\begin{column}{0.46\textwidth}
\begin{lstlisting}[language=R,basicstyle=\ttfamily\scriptsize]
# Install mmrm package
install.packages("mmrm")
library(mmrm)

# Load example data
data("fev_data", package = "mmrm")
\end{lstlisting}
\end{column}
\end{columns}

\vspace{0.1cm}
\begin{exampleblock}{\footnotesize fev\_data Structure}
\footnotesize
\begin{itemize}
    \item \texttt{FEV1}：结局变量（肺功能指标）
    \item \texttt{ARMCD}：治疗组（TRT / PBO）
    \item \texttt{AVISIT}：访视（VIS1--VIS4）
    \item \texttt{USUBJID}：受试者ID
\end{itemize}
\end{exampleblock}
\end{frame}

%---------------------------------------------------
\begin{frame}[fragile]{mmrm基本语法}
\begin{block}{标准MMRM模型}
\begin{lstlisting}[language=R,basicstyle=\ttfamily\scriptsize]
fit <- mmrm(
  formula = FEV1 ~ RACE + SEX + ARMCD * AVISIT 
           + us(AVISIT | USUBJID),
  data = fev_data, reml = TRUE
)
\end{lstlisting}
\end{block}

\textbf{公式解析：}
\begin{itemize}
    \item \texttt{FEV1 \textasciitilde\ ...}：结局变量
    \item \texttt{RACE + SEX}：协变量
    \item \texttt{ARMCD * AVISIT}：治疗、访视及其交互
    \item \texttt{us(AVISIT | USUBJID)}：\textbf{非结构化协方差}，按访视在受试者内
\end{itemize}

\vspace{0.1cm}
\begin{exampleblock}{协方差结构切换}
\texttt{us()} $\rightarrow$ \texttt{cs()}（复合对称）\quad
\texttt{us()} $\rightarrow$ \texttt{ar1()}（一阶自回归）\quad
\texttt{us()} $\rightarrow$ \texttt{toep()}（Toeplitz）
\end{exampleblock}
\end{frame}

%---------------------------------------------------
\begin{frame}[fragile]{自由度调整方法}
\begin{block}{mmrm支持的自由度方法}
\begin{lstlisting}[language=R,basicstyle=\ttfamily\scriptsize]
# Satterthwaite (default)
fit_sw <- mmrm(
  formula = FEV1 ~ ARMCD * AVISIT + us(AVISIT | USUBJID),
  data = fev_data, method = "Satterthwaite"
)

# Kenward-Roger
fit_kr <- mmrm(
  formula = FEV1 ~ ARMCD * AVISIT + us(AVISIT | USUBJID),
  data = fev_data, method = "Kenward-Roger"
)
\end{lstlisting}
\end{block}

\vspace{0.1cm}
\begin{columns}[T,totalwidth=\textwidth]
\begin{column}{0.48\textwidth}
\textbf{Satterthwaite}
\begin{itemize}
    \item 基于渐近方差
    \item 计算更快
    \item 大样本下与KR接近
\end{itemize}
\end{column}
\begin{column}{0.48\textwidth}
\textbf{Kenward-Roger}
\begin{itemize}
    \item 调整系数方差-协方差
    \item 小样本更保守
    \item 确证性试验常用
\end{itemize}
\end{column}
\end{columns}
\end{frame}

%---------------------------------------------------
\begin{frame}[fragile]{emmeans：最小二乘均值与对比}
\begin{columns}[T,totalwidth=\textwidth]
\begin{column}{0.48\textwidth}
\begin{block}{估计各访视的治疗差异}
\begin{lstlisting}[language=R,basicstyle=\ttfamily\scriptsize]
library(emmeans)

# LS means by visit
emm <- emmeans(fit, ~ ARMCD | AVISIT)

# Treatment contrast
pairs(emm, reverse = TRUE)
\end{lstlisting}
\end{block}
\end{column}
\begin{column}{0.48\textwidth}
\textbf{\footnotesize Output (partial):}
\begin{lstlisting}[language=R,numbers=none,basicstyle=\ttfamily\tiny]
AVISIT = VIS1:
 contrast estimate SE  df t.ratio p.value
 TRT-PBO   3.77 1.07 146  3.514 0.0006

AVISIT = VIS4:
 contrast estimate SE  df t.ratio p.value
 TRT-PBO   4.40 1.68 133  2.617 0.0099
\end{lstlisting}
\end{column}
\end{columns}

\vspace{0.1cm}
\begin{alertblock}{Primary Endpoint}
Treatment difference at final visit (VIS4) is usually the \textbf{primary endpoint estimate}.
\end{alertblock}
\end{frame}

%---------------------------------------------------
\begin{frame}[fragile]{mmrm结果解读}
\begin{block}{summary(fit) Output Highlights}
\begin{lstlisting}[language=R,numbers=none,basicstyle=\ttfamily\tiny]
Coefficients:
                      Estimate Std.Error    df t value Pr(>|t|)
(Intercept)           30.77748   0.88656 218.8  34.715  <2e-16 ***
ARMCDTRT               3.77423   1.07415 145.6   3.514 0.000589 ***
AVISITVIS4            15.05390   1.31281 138.5  11.467  <2e-16 ***
ARMCDTRT:AVISITVIS4    0.62423   1.85085 129.7   0.337 0.736463
\end{lstlisting}
\end{block}

\textbf{Interpretation Checklist:}
\begin{enumerate}
    \item \textbf{Convergence}: confirm optimizer converged (code = 0)
    \item \textbf{Covariance}: check variances positive, correlations reasonable
    \item \textbf{Fixed effects}: focus on treatment$\times$visit interaction
    \item \textbf{Final visit}: use emmeans to get LSM difference
    \item \textbf{Degrees of freedom}: note Satterthwaite/KR adjusted df
\end{enumerate}
\end{frame}

%---------------------------------------------------
\begin{frame}[fragile]{nlme::gls：替代实现}
\begin{block}{使用nlme包复现UN-MMRM}
\begin{lstlisting}[language=R,basicstyle=\ttfamily\scriptsize]
library(nlme)

fit_gls <- gls(
  FEV1 ~ RACE + SEX + ARMCD * AVISIT,
  data = fev_data,
  correlation = corSymm(form = ~ as.integer(AVISIT) | USUBJID),
  weights = varIdent(form = ~ 1 | AVISIT),
  na.action = na.omit, method = "REML"
)
\end{lstlisting}
\end{block}

\textbf{组件对应：}
\begin{itemize}
    \item \texttt{corSymm}：非结构化\textbf{相关}矩阵
    \item \texttt{varIdent}：访视特异性\textbf{方差}（异质性）
    \item 两者组合 = 非结构化协方差矩阵
\end{itemize}

\begin{alertblock}{注意事项}
\texttt{gls}较慢、更易收敛失败，且\textbf{不直接支持Kenward-Roger自由度}。新分析推荐使用\texttt{mmrm}包。
\end{alertblock}
\end{frame}

%===================================================
\section{SAS实现：PROC MIXED}
%===================================================

%---------------------------------------------------
\begin{frame}[fragile]{PROC MIXED基本语法}
\begin{columns}[T,totalwidth=\textwidth]
\begin{column}{0.52\textwidth}
\begin{block}{SAS代码：UN-MMRM + KR}
\begin{lstlisting}[language=R,numbers=none,basicstyle=\ttfamily\scriptsize]
proc mixed data=fev method=reml;
  class usubjid armcd avisit race sex;
  model fev1 = race sex armcd avisit 
         armcd*avisit / ddfm=kr;
  repeated avisit / subject=usubjid 
           type=un;
  lsmeans armcd*avisit / diff;
run;
\end{lstlisting}
\end{block}
\end{column}
\begin{column}{0.44\textwidth}
\textbf{关键语句解析：}
\begin{itemize}
    \item \texttt{class}：声明分类变量
    \item \texttt{model}：固定效应模型
    \item \texttt{ddfm=kr}：Kenward-Roger
    \item \texttt{repeated}：R-side协方差
    \item \texttt{type=un}：非结构化
    \item \texttt{lsmeans}：最小二乘均值
\end{itemize}
\end{column}
\end{columns}
\end{frame}

%---------------------------------------------------
\begin{frame}{SAS关键选项详解}
\begin{columns}[T,totalwidth=\textwidth]
\begin{column}{0.48\textwidth}
\begin{block}{MODEL语句选项}
\begin{itemize}
    \item \texttt{ddfm=kr}：Kenward-Roger
    \item \texttt{ddfm=kr(linear)}：线性KR
    \item \texttt{ddfm=satterthwaite}：Satterthwaite
    \item \texttt{ddfm=bw}：Between-Within
    \item \texttt{ddfm=residual}：残差df
    \item \texttt{solution}：输出固定效应解
\end{itemize}
\end{block}
\end{column}
\begin{column}{0.48\textwidth}
\begin{block}{REPEATED语句选项}
\begin{itemize}
    \item \texttt{type=un}：非结构化
    \item \texttt{type=cs}：复合对称
    \item \texttt{type=ar(1)}：一阶自回归
    \item \texttt{type=toep}：Toeplitz
    \item \texttt{type=unh}：异质性UN
    \item \texttt{group=}：分组协方差
\end{itemize}
\end{block}
\end{column}
\end{columns}

\vspace{0.2cm}
\begin{exampleblock}{R与SAS的对应关系}
\centering
\texttt{us(AVISIT | USUBJID)} $\leftrightarrow$ \texttt{repeated AVISIT / subject=USUBJID type=UN}

\texttt{method="Kenward-Roger"} $\leftrightarrow$ \texttt{ddfm=KR}
\end{exampleblock}
\end{frame}

%---------------------------------------------------
\begin{frame}{SAS输出解读}
\begin{block}{主要输出表格}
\begin{enumerate}
    \item \textbf{模型信息}：方法（REML/ML）、协方差结构、受试者数
    \item \textbf{收敛状态}：确认迭代收敛
    \item \textbf{协方差参数估计}：UN结构的方差和相关
    \item \textbf{拟合标准}：AIC、BIC、-2Res Log Likelihood
    \item \textbf{固定效应解}：系数估计、标准误、df、t值、p值
    \item \textbf{LSM差异}：各访视的治疗组对比
\end{enumerate}
\end{block}

\vspace{0.2cm}
\begin{alertblock}{与R结果的一致性}
\begin{itemize}
    \item 使用相同数据和相同结构时，SAS与R的\textbf{系数估计应一致}
    \item R mmrm的\texttt{vcov="Kenward-Roger-Linear"}可匹配SAS的线性KR
\end{itemize}
\end{alertblock}
\end{frame}

%===================================================
\section{R与SAS代码对照与案例演示}
%===================================================

%---------------------------------------------------
\begin{frame}{R与SAS代码对照表}
\begin{center}
\small
\begin{tabular}{p{5.5cm}p{5.5cm}}
\toprule
\textbf{R (mmrm)} & \textbf{SAS (PROC MIXED)} \\
\midrule
\texttt{mmrm(formula=..., data=...)} & \texttt{proc mixed data=...;} \\
\texttt{FEV1 \textasciitilde\ ARMCD * AVISIT} & \texttt{model fev1 = armcd avisit armcd*avisit;} \\
\texttt{+ RACE + SEX} & \texttt{+ race sex} \\
\texttt{us(AVISIT | USUBJID)} & \texttt{repeated avisit / subject=usubjid type=un;} \\
\texttt{reml = TRUE} & \texttt{method=reml} \\
\texttt{method="Kenward-Roger"} & \texttt{ddfm=kr} \\
\texttt{emmeans(fit, \textasciitilde\ ARMCD | AVISIT)} & \texttt{lsmeans armcd*avisit;} \\
\texttt{pairs(..., reverse=TRUE)} & \texttt{lsmeans ... / diff;} \\
\bottomrule
\end{tabular}
\end{center}

\vspace{0.2cm}
\begin{exampleblock}{关键差异}
\begin{itemize}
    \item R使用公式语法，SAS使用语句式语法
    \item R的\texttt{us()}自动处理异质性，SAS的\texttt{TYPE=UN}同理
\end{itemize}
\end{exampleblock}
\end{frame}

%---------------------------------------------------
\begin{frame}{完整案例：FEV1肺功能试验}
\begin{block}{研究背景}
随机对照试验，比较治疗组（TRT）与安慰剂组（PBO）对FEV1（肺功能）的影响。测量时间点：VIS1--VIS4。
\end{block}

\textbf{分析目标：}
\begin{enumerate}
    \item 估计各访视两组间的FEV1差异
    \item 检验最终访视（VIS4）的治疗效应
    \item 在MAR假设下处理缺失数据
\end{enumerate}

\vspace{0.1cm}
\begin{alertblock}{模型预设}
\begin{itemize}
    \item 协方差结构：UN（主要），Toeplitz（fallback）
    \item 自由度：Kenward-Roger\quad 估计方法：REML\quad 主要终点：VIS4的LSM差异
\end{itemize}
\end{alertblock}
\end{frame}

%---------------------------------------------------
\begin{frame}[fragile]{案例：R完整代码}
\begin{lstlisting}[language=R,basicstyle=\ttfamily\scriptsize]
library(mmrm)
library(emmeans)

# 1. Fit UN-MMRM
fit <- mmrm(
  formula = FEV1 ~ RACE + SEX + ARMCD * AVISIT + us(AVISIT | USUBJID),
  data = fev_data, method = "Kenward-Roger"
)

# 2. Check convergence
component(fit, "convergence")  # should be 0

# 3. View covariance estimate
VarCorr(fit)

# 4. LS means and contrasts
emm <- emmeans(fit, ~ ARMCD | AVISIT)
contrasts <- pairs(emm, reverse = TRUE)

# 5. Extract final visit results
summary(contrasts)  # focus on VIS4
\end{lstlisting}
\end{frame}

%---------------------------------------------------
\begin{frame}[fragile]{案例：SAS完整代码}
\begin{lstlisting}[language=R,numbers=none,basicstyle=\ttfamily\scriptsize]
proc mixed data=fev method=reml;
  class usubjid armcd avisit race sex;
  model fev1 = race sex armcd avisit armcd*avisit / ddfm=kr solution;
  repeated avisit / subject=usubjid type=un r rcorr;
  lsmeans armcd*avisit / diff slice=avisit;
  estimate 'TRT-PBO at VIS4' armcd 1 -1 armcd*avisit 0 0 0 1  0 0 0 -1;
run;
\end{lstlisting}

\vspace{0.2cm}
\begin{exampleblock}{SAS-specific Options}
\begin{itemize}
    \item \texttt{r} / \texttt{rcorr}: output estimated covariance/correlation matrix
    \item \texttt{slice=AVISIT}: stratified test by visit
    \item \texttt{estimate}: custom linear contrast
\end{itemize}
\end{exampleblock}
\end{frame}

%---------------------------------------------------
\begin{frame}{结果一致性验证}
\begin{block}{验证要点}
当R与SAS使用\textbf{相同模型}时，应检查：
\begin{enumerate}
    \item 固定效应系数估计（小数点后3-4位一致）
    \item 协方差参数估计
    \item 最终访视的LSM差异及其标准误
    \item 模型拟合统计量（-2 log likelihood）
\end{enumerate}
\end{block}

\vspace{0.1cm}
\begin{center}
\small
\begin{tabular}{lcc}
\toprule
\textbf{指标} & \textbf{R mmrm} & \textbf{SAS PROC MIXED} \\
\midrule
VIS4 LSM差异 & 4.40 & 4.40 \\
标准误 & 1.68 & 1.68 \\
df (KR) & 133 & 133 \\
t值 & 2.617 & 2.617 \\
p值 & 0.0099 & 0.0099 \\
-2 Res Log L & 3387.4 & 3387.4 \\
\bottomrule
\end{tabular}
\end{center}

\vspace{0.1cm}
\begin{alertblock}{注意}
若使用不同的KR实现（线性vs非线性），标准误可能有微小差异。R中设置\texttt{vcov="Kenward-Roger-Linear"}可匹配SAS默认行为。
\end{alertblock}
\end{frame}

%===================================================
\section{高级主题与注意事项}
%===================================================

%---------------------------------------------------
\begin{frame}{Estimand框架与MMRM}
\begin{block}{ICH E9(R1) addendum}
MMRM是\textbf{估计方法}，不是estimand。在确证性试验中：
\begin{itemize}
    \item 先定义estimand（目标估计量）
    \item 再选择匹配的分析方法
\end{itemize}
\end{block}

\textbf{MMRM与策略的匹配：}
\begin{itemize}
    \item \textbf{假设策略}（Hypothetical）：MMRM + MAR最自然
    \item \textbf{治疗政策策略}（Treatment-policy）：可能需要不同缺失处理
    \item \textbf{复合策略}（Composite）：MMRM可能不适用
\end{itemize}

\begin{alertblock}{关键原则}
MAR假设必须是有意匹配所选estimand的\textbf{预设选择}，而非事后想法。若怀疑MNAR，需预设敏感性分析。
\end{alertblock}
\end{frame}

%---------------------------------------------------
\begin{frame}[fragile]{收敛问题与诊断}
\begin{block}{Common Convergence Issues}
\textbf{1. UN does not converge:} many visits, small sample, complex missing patterns

\textbf{2. Covariance not positive definite:} negative variance or $|\rho|>1$

\textbf{3. Hessian not positive definite:} flat optimization surface, non-identifiable parameters
\end{block}

\vspace{0.1cm}
\textbf{Solutions:}

\textbullet~Preset fallback covariance structure (Toeplitz, AR1)

\textbullet~Try different optimizers (L-BFGS-B, BFGS, nlminb)

\textbullet~Check data: outliers, visit coding, missing patterns

\textbullet~Simplify model: reduce interactions, merge sparse categories

\vspace{0.1cm}
\begin{lstlisting}[language=R,basicstyle=\ttfamily\scriptsize]
# Try multiple optimizers
fit <- mmrm(..., control = mmrm_control(
  optimizer = c("L-BFGS-B", "BFGS", "nlminb")
))
\end{lstlisting}
\end{frame}

%---------------------------------------------------
\begin{frame}[fragile]{模型诊断与残差分析}
\begin{block}{mmrm支持的残差类型}
\begin{lstlisting}[language=R,basicstyle=\ttfamily\scriptsize]
# Raw residuals
res_resp <- residuals(fit, type = "response")

# Pearson residuals (recommended for diagnostics)
res_pearson <- residuals(fit, type = "pearson")

# Standardized residuals (approx N(0,1), for normality check)
res_norm <- residuals(fit, type = "normalized")
\end{lstlisting}
\end{block}

\textbf{诊断清单：}
\begin{enumerate}
    \item \textbf{正态性}：标准化残差的Q-Q图
    \item \textbf{异常值}：Pearson残差的绝对值 $> 3$
    \item \textbf{协方差结构}：比较不同结构的AIC/BIC
    \item \textbf{影响点}：检查极端值对估计的影响
\end{enumerate}

\begin{alertblock}{注意}
MMRM允许异方差，因此\textbf{原始残差}应谨慎解读。推荐使用\textbf{Pearson残差}或\textbf{标准化残差}进行诊断。
\end{alertblock}
\end{frame}

%===================================================
\section{总结}
%===================================================

%---------------------------------------------------
\begin{frame}{要点总结}
\begin{block}{核心要点}
\begin{enumerate}
    \item \textbf{MMRM是LMM的特例}：省略G-side随机效应，仅通过R-side结构化协方差建模within-subject相关
    \item \textbf{确证性试验金标准}：在MAR下无需插补，利用所有观测数据
    \item \textbf{UN是默认结构}：假设最少，但需预设fallback
    \item \textbf{访视作分类变量}：允许治疗效应随时间呈任意形状
    \item \textbf{R与SAS可互验}：相同模型应产生一致结果
\end{enumerate}
\end{block}

\vspace{0.2cm}
\begin{columns}[T,totalwidth=\textwidth]
\begin{column}{0.48\textwidth}
\textbf{选择LMM当：}
\begin{itemize}
    \item 关注受试者间变异
    \item 多层级聚类结构
    \item 个体预测有意义
\end{itemize}
\end{column}
\begin{column}{0.48\textwidth}
\textbf{选择MMRM当：}
\begin{itemize}
    \item 组间比较为主要目标
    \item 计划访视的纵向数据
    \item 需处理MAR缺失
\end{itemize}
\end{column}
\end{columns}
\end{frame}

%---------------------------------------------------
\begin{frame}{参考文献}
\small
\begin{enumerate}
    \item Laird NM, Ware JH. Random-effects models for longitudinal data. \textit{Biometrics}. 1982;38(4):963-974.
    \item Pinheiro J, Bates DM. \textit{Mixed-Effects Models in S and S-PLUS}. Springer; 2000.
    \item Mallinckrodt CH, et al. Recommendations for the Primary Analysis of Continuous Endpoints in Longitudinal Clinical Trials. \textit{Drug Information Journal}. 2008;42(4):303-319.
    \item Mallinckrodt C, Lipkovich I. \textit{Analyzing Longitudinal Clinical Trial Data: A Practical Guide}. Chapman \& Hall/CRC; 2017.
    \item National Research Council. \textit{The Prevention and Treatment of Missing Data in Clinical Trials}. The National Academies Press; 2010.
    \item ICH E9(R1). Addendum on Estimands and Sensitivity Analysis in Clinical Trials. 2019.
    \item Cnaan A, Laird NM, Slasor P. Using the general linear mixed model to analyse unbalanced repeated measures. \textit{Statistics in Medicine}. 1997;16(20):2349-2380.
    \item Littell RC, et al. Modelling covariance structure in the analysis of repeated measures data. \textit{Statistics in Medicine}. 2000;19(13):1793-1819.
    \item openpharma/mmrm R package documentation. \url{https://openpharma.github.io/mmrm/}
    \item SAS/STAT 9.2 User's Guide: PROC MIXED. SAS Institute Inc.
\end{enumerate}
\end{frame}

%---------------------------------------------------
\begin{frame}[plain]
\begin{center}
\vspace{2cm}
{\Huge \textcolor{darkblue}{谢谢！}}

\vspace{1cm}
{\Large Questions \& Discussion}



\end{center}
\end{frame}

\end{document}
